% =====================================================================
% Relativistic Normal Modes and Exchange of Stabilities
% Chapter II, Sections 10--12: Relativistic Modification
% =====================================================================
% Requires: rel_preamble.tex (loaded by rel_main.tex)
%
% Classical reference: Chandrasekhar, Ch. II, §§10--12
% Relativistic framework: Israel-Stewart causal dissipation theory
% =====================================================================

\chapter*{Chapter II (Relativistic): The B\'enard Problem --- Normal Mode
  Analysis, Exchange of Stabilities, and the Marginal State}
\addcontentsline{toc}{chapter}{Chapter II (Rel.): Normal Modes \&
  Exchange of Stabilities}

%----------------------------------------------------------------------
\section{The analysis into normal modes: relativistic}
\label{sec:rel-10}
%----------------------------------------------------------------------

\subsection{Covariant perturbation ansatz}
\label{sec:rel-10a}

In the classical theory (\S~10 of Chapter~II) every perturbed quantity is
expanded with the factor $\exp[i(k_x x + k_y y) + pt]$, where $p$ is an
unrestricted complex growth rate and $k = (k_x^2+k_y^2)^{1/2}$ is the
horizontal wave number.  The Newtonian dispersion relation that follows from
the linearised Boussinesq equations is a cubic polynomial in $p$, and all
three roots are physical.

In the relativistic setting the situation changes in two essential respects:

\begin{enumerate}
\item \textbf{Causal propagation.}  The perturbation equations must be
  derived from the Israel--Stewart theory, which replaces the parabolic
  (instantaneous-diffusion) heat and viscous-stress equations with hyperbolic
  relaxation equations.  Every signal speed that emerges from the dispersion
  relation must be bounded by~$c$.

\item \textbf{Additional modes.}  Each Israel--Stewart relaxation equation
  adds an independent dynamical degree of freedom.  The dispersion relation
  consequently becomes a \emph{higher-order} polynomial in the growth
  parameter, producing extra roots that represent transient relaxation modes
  with no Newtonian counterpart.
\end{enumerate}

We work in the local rest frame of the static equilibrium and retain factors
of $c$ explicitly throughout, so that the non-relativistic limit $c\to\infty$
can be taken at every stage.

\medskip
\noindent\textbf{Perturbation ansatz.}
Ascribe to every linearised quantity (the vertical velocity perturbation~$W$,
the temperature perturbation~$\Theta$, the vertical vorticity
perturbation~$Z$, the heat-flux perturbation~$\delta q^{\mu}$, and the
shear-stress perturbation~$\delta\pi^{\mu\nu}$) a dependence on the
horizontal coordinates and time of the form
\begin{equation}
\label{eq:rel-2-88}
\exp\!\bigl[i(k_x\, x + k_y\, y) + \sigma t\bigr],
\end{equation}
where $\sigma$ may be complex and $k = (k_x^2+k_y^2)^{1/2}$.  The vertical
dependence remains a function to be determined: $W(z)$, $\Theta(z)$, etc.

Introducing the non-dimensional wave number $a = kd$ and the non-dimensional
growth rate $\sigma_* = \sigma d^2/\nu$ (as in the classical \S~10), together
with the operators
\begin{equation}
\label{eq:rel-2-91}
D \equiv \frac{d}{dz}, \qquad
\mathcal{D}^2 \equiv D^2 - a^2,
\end{equation}
the linearised Israel--Stewart equations for the vertical modes reduce to the
following system.

\subsection{Relativistic normal-mode equations}
\label{sec:rel-10b}

The \emph{momentum equation} for the vertical velocity amplitude reads
\begin{equation}
\label{eq:rel-2-92}
\Bigl[\mathcal{D}^2(\mathcal{D}^2 - \sigma_*)
  + \frac{\sigma_* \taupi \nu}{d^2}\,\mathcal{D}^2
      (\mathcal{D}^2 - \sigma_* - d^2/(\taupi\nu))\Bigr] W
= -\Rarel\, a^2 \Theta,
\end{equation}
where $\Rarel$ is the relativistic Rayleigh number
\begin{equation}
\label{eq:rel-Rayleigh}
\Rarel = \frac{g\alpha\beta\, d^4}{\kappa\nu}
  \left[1 + \order{\frac{v^2}{c^2}}\right],
\end{equation}
and $\taupi$ is the shear-stress relaxation time of the Israel--Stewart
theory.  The first line of~\eqref{eq:rel-2-92} is exactly Chandrasekhar's
equation~(92); the second line is the relativistic (Israel--Stewart)
correction, which vanishes in the limit $\taupi\to 0$.

The \emph{energy equation} for the temperature amplitude, incorporating
the causal heat-flux relaxation time~$\tauq$, becomes
\begin{equation}
\label{eq:rel-2-93}
\Bigl[(\mathcal{D}^2 - \mathrm{p}\,\sigma_*)
  + \tauq\,\frac{\kappa}{d^2}\,\sigma_*
    \bigl(\mathcal{D}^2 - \mathrm{p}\,\sigma_*
      - d^2/(\tauq\kappa)\bigr)\Bigr]\Theta
= -\frac{\beta d^2}{\kappa}\, W,
\end{equation}
where $\mathrm{p} = \nu/\kappa$ is the Prandtl number.  Again the classical
equation~(93) is recovered when $\tauq\to 0$.

The \emph{vorticity equation} acquires an analogous relaxation correction:
\begin{equation}
\label{eq:rel-2-94}
\bigl[\mathcal{D}^2 - \sigma_*
  + \taupi(\nu/d^2)\,\sigma_*(\mathcal{D}^2 - \sigma_*
    - d^2/(\taupi\nu))\bigr] Z = 0.
\end{equation}
The boundary conditions remain unchanged from the classical case:
\begin{equation}
\label{eq:rel-2-95}
\Theta = 0, \quad W = 0 \qquad\text{for } z = 0 \text{ and } 1,
\end{equation}
together with either $DW = 0$ (rigid) or $D^2 W = 0$ (free) on each
bounding surface.

\subsection{Dispersion relation and mode identification}
\label{sec:rel-10c}

Eliminating $\Theta$ between~\eqref{eq:rel-2-92} and~\eqref{eq:rel-2-93}
and seeking solutions proportional to $\sin(n\pi z)$ (the ``both free''
boundary prototype), one obtains the dispersion relation
\begin{equation}
\label{eq:rel-dispersion}
\mathcal{P}(\sigma_*;\, a^2,\, \Rarel,\, \taupi,\, \tauq) = 0,
\end{equation}
where $\mathcal{P}$ is a polynomial of degree \emph{five} in~$\sigma_*$
(compared to the cubic of the classical theory).

The five roots decompose as follows:
\begin{enumerate}
\item[\textbf{(i)}] \textbf{Three physical (hydrodynamic) modes} ---
  continuously connected to the three roots of the classical dispersion
  relation as $\taupi, \tauq \to 0$.  These govern the onset of thermal
  convection.

\item[\textbf{(ii)}] \textbf{Two transient (relaxation) modes} --- with
  $\operatorname{Re}(\sigma_*) \sim -d^2/(\taupi\nu)$ or
  $\sim -d^2/(\tauq\kappa)$ for small relaxation times.  These modes decay
  rapidly on the relaxation time-scale and do not participate in the
  instability.  They are an artefact of the extended thermodynamic theory
  and have no classical counterpart.
\end{enumerate}

\begin{relcorrection}
In the Newtonian theory the dispersion relation is cubic in~$\sigma_*$
and all three modes are physical.  The Israel--Stewart extension raises the
polynomial order by two, adding one mode per relaxation channel (heat and
shear).  These ``spurious'' (from the Newtonian viewpoint) roots are
essential for relativistic causality: they ensure that signals propagate at
finite speed.
\end{relcorrection}

\subsection{Causality constraint on group velocities}
\label{sec:rel-10d}

\begin{causalitycheck}
For any mode with dispersion relation $\sigma_*(a)$, define the group
velocity
\begin{equation}
\label{eq:rel-vg}
v_g = \left|\frac{\partial\,\operatorname{Re}(\sigma_*)}
            {\partial a}\right| \frac{\nu}{d}.
\end{equation}
Causality requires $v_g \le c$ for every branch.  This is guaranteed
provided the relaxation times satisfy
\begin{equation}
\label{eq:rel-causality-shear}
\taupi \ge \frac{\shearvisc}{\enthalpy\, c^2}\,,
\qquad
\tauq \ge \frac{\thermcond\, T}{\enthalpy\, c^2}\,,
\end{equation}
which are precisely the Israel--Stewart causality bounds.  When these hold,
the maximum signal speed of the perturbation system is bounded by
$\sqrt{\nu/\taupi}$ and $\sqrt{\kappa/\tauq}$ (in the respective channels),
both of which are $\le c$ by~\eqref{eq:rel-causality-shear}.

For the three physical (hydrodynamic) modes, which govern convective
onset, the group velocities are $\order{v}$ with $v \ll c$; causality
is satisfied trivially.  The constraint is non-trivial only for the
transient relaxation modes at large~$a$ (short wavelengths).
\end{causalitycheck}

\subsection{Horizontal velocity components}
\label{sec:rel-10e}

The reconstruction of the horizontal velocities $u$, $v$ from $W$ and $Z$ via
the stream-function decomposition (equations~(106)--(111) of \S~10a) is
purely kinematic and carries over unchanged to the relativistic case,
provided one works in the local rest frame.  No additional relativistic
corrections appear at this step.

%----------------------------------------------------------------------
\section{The principle of the exchange of stabilities: relativistic}
\label{sec:rel-11}
%----------------------------------------------------------------------

\subsection{Statement of the problem}
\label{sec:rel-11a}

The classical proof (\S~11) demonstrates that for $R > 0$ the growth
rate $\sigma$ is purely real, so the transition from stability to instability
occurs through a \emph{stationary} (non-oscillatory) marginal state.  The key
question in the relativistic extension is:

\begin{quote}
\emph{Does $\sigma_*$ remain real at the marginal state when the
Israel--Stewart relaxation terms are included?}
\end{quote}

We shall prove that the answer is \emph{yes} for the three physical
modes, up to corrections that vanish when the relaxation times tend to
zero.  The two transient modes are always heavily damped and do not affect
the exchange-of-stabilities argument.

\subsection{Proof for the physical modes}
\label{sec:rel-11b}

Define, as in the classical analysis,
\begin{equation}
\label{eq:rel-2-112}
G = \mathcal{D}^2 W, \qquad
F = \mathcal{D}^2(\mathcal{D}^2 - \sigma_*) W
  = (\mathcal{D}^2 - \sigma_*) G.
\end{equation}
From the relativistic normal-mode equations~\eqref{eq:rel-2-92}
and~\eqref{eq:rel-2-93}, after eliminating $\Theta$, we obtain
\begin{equation}
\label{eq:rel-2-115}
(\mathcal{D}^2 - \mathrm{p}\,\sigma_*) F
  + \mathcal{R}_{\mathrm{IS}}[\sigma_*, F, W]
  = -\Rarel\, a^2 W,
\end{equation}
where
\begin{equation}
\label{eq:rel-IS-remainder}
\mathcal{R}_{\mathrm{IS}}
= \frac{\taupi\nu}{d^2}\,\sigma_*\,
    \bigl[\mathcal{D}^2 F
      - (\sigma_* + d^2/(\taupi\nu))\, F\bigr]
  + \frac{\tauq\kappa}{d^2}\,\sigma_*\,
    \bigl[\mathcal{D}^2 - \mathrm{p}\,\sigma_*
      - d^2/(\tauq\kappa)\bigr]
    \bigl[\mathcal{D}^2(\mathcal{D}^2 - \sigma_*) W\bigr]
\end{equation}
collects all Israel--Stewart corrections, and
$\mathcal{R}_{\mathrm{IS}} \to 0$ as $\taupi, \tauq \to 0$.

Multiply~\eqref{eq:rel-2-115} by $F^*$ and integrate over
$z \in [0,1]$.  Following Chandrasekhar's integration-by-parts strategy
exactly (boundary terms vanish by the same argument as in the classical
case), we arrive at
\begin{equation}
\label{eq:rel-2-118}
\int_0^1\!\bigl\{|DF|^2 + a^2|F|^2
  + \mathrm{p}\,\sigma_*|F|^2\bigr\}\,dz
  + \int_0^1 F^*\,\mathcal{R}_{\mathrm{IS}}\,dz
  = \Rarel\, a^2\!\int_0^1 W F^*\,dz.
\end{equation}
The right-hand side is treated identically to the classical
equations~(119)--(123), yielding
\begin{equation}
\label{eq:rel-2-124}
\int_0^1\!\bigl\{|DF|^2 + a^2|F|^2
  + \mathrm{p}\,\sigma_*|F|^2\bigr\}\,dz
  + \int_0^1 F^*\,\mathcal{R}_{\mathrm{IS}}\,dz
  = \Rarel\, a^2\!\int_0^1\!\bigl\{|G|^2
    + \sigma_*^*\bigl[|DW|^2 + a^2|W|^2\bigr]\bigr\}\,dz.
\end{equation}
Taking the imaginary part:
\begin{equation}
\label{eq:rel-2-125}
\operatorname{Im}(\sigma_*)
  \left\{\mathrm{p}\!\int_0^1|F|^2\,dz
    + \Rarel\, a^2\!\int_0^1
      \bigl[|DW|^2 + a^2|W|^2\bigr]\,dz\right\}
= -\operatorname{Im}\!\int_0^1 F^*\,\mathcal{R}_{\mathrm{IS}}\,dz.
\end{equation}
The left-hand factor in braces is positive definite for $\Rarel > 0$.
Therefore:
\begin{equation}
\label{eq:rel-2-126}
\operatorname{Im}(\sigma_*) = \order{\taupi,\, \tauq}.
\end{equation}

\begin{relcorrection}
\textbf{Exchange of stabilities (relativistic).}
In the classical theory ($\taupi = \tauq = 0$) the right-hand side
of~\eqref{eq:rel-2-125} vanishes identically and one concludes
$\operatorname{Im}(\sigma) = 0$ exactly.

With Israel--Stewart corrections, $\operatorname{Im}(\sigma_*)$ receives a
contribution of order $\max(\taupi\nu/d^2,\, \tauq\kappa/d^2)$, which is
the ratio of the \emph{microscopic} relaxation time to the
\emph{macroscopic} viscous diffusion time.  For any laboratory or
astrophysical fluid this ratio is exceedingly small ($\sim 10^{-10}$ or
less), so the principle of exchange of stabilities holds to extraordinary
accuracy.

Formally, the principle is \emph{exact} in the non-relativistic limit
$c \to \infty$ (equivalently $\taupi, \tauq \to 0$), and receives
corrections suppressed by $v^2/c^2$:
\begin{equation}
\label{eq:rel-exchange-correction}
\operatorname{Im}(\sigma_*) = \order{\frac{v^2}{c^2}}.
\end{equation}
\end{relcorrection}

\subsection{The transient relaxation modes}
\label{sec:rel-11c}

The two extra roots of the fifth-order dispersion
relation~\eqref{eq:rel-dispersion} have
\begin{equation}
\label{eq:rel-transient}
\operatorname{Re}(\sigma_*^{(\mathrm{trans})})
\approx -\frac{d^2}{\taupi\nu}
\quad\text{or}\quad
-\frac{d^2}{\tauq\kappa}.
\end{equation}
Since $d^2/(\taupi\nu) \gg 1$ and $d^2/(\tauq\kappa) \gg 1$ for any
macroscopic system, these modes are strongly damped.  They never become
marginal and therefore play no role in the onset of convection.  Their
presence is required solely by the hyperbolic character of the
Israel--Stewart equations and ensures that no perturbation signal propagates
faster than~$c$.

%----------------------------------------------------------------------
\section{The marginal state and the relativistic characteristic value
  problem}
\label{sec:rel-12}
%----------------------------------------------------------------------

\subsection{Marginal-state equations}
\label{sec:rel-12a}

Since the exchange of stabilities is valid (to the accuracy established
in \S~\ref{sec:rel-11}), the marginal state is stationary:
$\sigma_* = 0$.  Setting $\sigma_* = 0$ in the normal-mode
equations~\eqref{eq:rel-2-92} and~\eqref{eq:rel-2-93} and observing that
every Israel--Stewart correction term carries a factor of $\sigma_*$
(hence vanishes at $\sigma_* = 0$), we obtain
\begin{equation}
\label{eq:rel-2-127}
(\mathcal{D}^2)^2 W
= (D^2-a^2)^2 W
= +\Rarel\, a^2 \Theta,
\end{equation}
\begin{equation}
\label{eq:rel-2-128}
(D^2-a^2)\Theta = -\frac{\beta d^2}{\kappa}\, W.
\end{equation}

\begin{relcorrection}
\textbf{Key observation.}  At marginal stability ($\sigma_* = 0$),
every Israel--Stewart relaxation term vanishes because each such term is
proportional to~$\sigma_*$.  The marginal-state equations therefore take
\emph{exactly the same form} as the classical equations~(127)--(128) of
Chandrasekhar, with $R$ replaced by $\Rarel$.

The entire effect of relativistic corrections is thus absorbed into the
definition of $\Rarel$.  This is physically sensible: the marginal state
is a time-independent configuration, so the hyperbolic (causal) modifications
to the time-evolution equations are invisible.
\end{relcorrection}

Eliminating $\Theta$:
\begin{equation}
\label{eq:rel-2-129}
(D^2 - a^2)^3 W = -\Rarel\, a^2\, W,
\end{equation}
with boundary conditions
\begin{equation}
\label{eq:rel-2-130}
W = 0, \quad (D^2-a^2)^2 W = 0 \qquad\text{for } z = 0 \text{ and } 1,
\end{equation}
and either $DW = 0$ (rigid) or $D^2 W = 0$ (free) on each surface.

Alternatively, eliminating $W$:
\begin{equation}
\label{eq:rel-2-131}
(D^2-a^2)^3\Theta = -\Rarel\, a^2\, \Theta,
\end{equation}
with boundary conditions
\begin{equation}
\label{eq:rel-2-132}
\Theta = 0, \quad (D^2-a^2)\Theta = 0 \qquad\text{for } z = 0 \text{ and } 1,
\end{equation}
and appropriate conditions on $D(D^2-a^2)\Theta$ depending on whether
each surface is rigid or free.

\subsection{The characteristic value problem for $\Rarel$}
\label{sec:rel-12b}

Equations~\eqref{eq:rel-2-129}--\eqref{eq:rel-2-130} (or equivalently
\eqref{eq:rel-2-131}--\eqref{eq:rel-2-132}) constitute a sixth-order
ordinary differential equation with six homogeneous boundary conditions.
A non-trivial solution exists only for particular values of $\Rarel$ at
given~$a^2$; we thus have a \emph{characteristic value problem} for the
relativistic Rayleigh number.

Let $\Rarel^{(c)}(a^2)$ denote the lowest eigenvalue for a given~$a^2$.
The critical relativistic Rayleigh number is
\begin{equation}
\label{eq:rel-Rac}
\Rarel^{(\mathrm{crit})}
= \min_{a^2} \Rarel^{(c)}(a^2).
\end{equation}

Since the eigenvalue problem at $\sigma_* = 0$ has exactly the same
structure as in the classical theory, we can write
\begin{equation}
\label{eq:rel-Rac-expansion}
\Rarel^{(\mathrm{crit})}
= R_c^{(\mathrm{class})}
  \left[1 + \order{\frac{v^2}{c^2}}\right],
\end{equation}
where $R_c^{(\mathrm{class})}$ is the classical critical Rayleigh number
(e.g.\ $R_c = 27\pi^4/4$ for both surfaces free).  The relativistic
correction enters solely through the $\order{v^2/c^2}$ terms in the
definition of $\Rarel$ (equation~\eqref{eq:rel-Rayleigh}).

The critical wave number $a_c$ at which $\Rarel^{(\mathrm{crit})}$ is
attained is likewise corrected only at $\order{v^2/c^2}$:
\begin{equation}
\label{eq:rel-ac}
a_c^{(\mathrm{rel})}
= a_c^{(\mathrm{class})}
  \left[1 + \order{\frac{v^2}{c^2}}\right].
\end{equation}

\subsection{Role of the relaxation modes in the eigenvalue problem}
\label{sec:rel-12c}

Although the transient relaxation modes (identified in
\S~\ref{sec:rel-10c}) do not participate in the instability, their presence
modifies the \emph{approach} to the marginal state.  Specifically, for
$\sigma_*$ small but non-zero, the dispersion relation receives
corrections of order $\sigma_* \taupi\nu/d^2$ and
$\sigma_* \tauq\kappa/d^2$.  These produce $\order{v^2/c^2}$ shifts in the
growth rate near onset, but as shown above, the marginal eigenvalue problem
itself is unaffected.

\begin{causalitycheck}
\textbf{Summary of causality verification for the eigenvalue problem.}

\begin{enumerate}
\item The three physical modes satisfy
  $|v_g| \ll c$ for all wave numbers relevant to convective onset.

\item The two transient modes have group velocities bounded by
  $\sqrt{\nu/\taupi}$ and $\sqrt{\kappa/\tauq}$, both $\le c$ by the
  Israel--Stewart causality bounds~\eqref{eq:rel-causality-shear}.

\item The marginal-state eigenvalue problem ($\sigma_* = 0$) involves no
  time derivatives and is automatically causal.

\item The critical Rayleigh number and critical wave number receive only
  $\order{v^2/c^2}$ corrections relative to their classical values.
\end{enumerate}

All mode frequencies therefore give subluminal group velocities, and the
relativistic eigenvalue problem is well-posed.
\end{causalitycheck}

\subsection{Non-relativistic limit}
\label{sec:rel-12d}

Taking $c \to \infty$ (hence $\taupi, \tauq \to 0$ and
$\Rarel \to R$), the entire system of equations in \S\S~\ref{sec:rel-10}--\ref{sec:rel-12} reduces to Chandrasekhar's classical
equations~(88)--(132) of Chapter~II, \S\S~10--12.  The five-mode dispersion
relation collapses to the classical cubic as the two transient modes acquire
$|\operatorname{Re}(\sigma_*)| \to \infty$ and decouple.  The exchange of
stabilities becomes exact, and the characteristic value problem for $R$
is recovered without modification.
